\documentclass[12pt]{article}

% Core math
\usepackage{amsmath, amssymb, amsthm}

% Diagrams
\usepackage{tikz-cd}
\usepackage{tikz}
\usetikzlibrary{decorations.markings, arrows.meta, positioning}

% References
\usepackage{hyperref}
\usepackage{cleveref}

% Graphics
\usepackage{graphicx}

% Page layout
\usepackage[margin=1in]{geometry}

% GrokRxiv sidebar
\usepackage{everypage}
\usepackage{xcolor}

% Listings for code blocks
\usepackage{listings}
\lstset{
  basicstyle=\ttfamily\footnotesize,
  columns=fullflexible,
  breaklines=true,
  showstringspaces=false,
  frame=single,
  framerule=0.4pt,
  rulecolor=\color{black!40},
  xleftmargin=0pt,
  xrightmargin=0pt
}

% Theorem environments
\newtheorem{theorem}{Theorem}[section]
\newtheorem{proposition}[theorem]{Proposition}
\newtheorem{lemma}[theorem]{Lemma}
\newtheorem{corollary}[theorem]{Corollary}
\theoremstyle{definition}
\newtheorem{definition}[theorem]{Definition}
\newtheorem{example}[theorem]{Example}
\theoremstyle{remark}
\newtheorem{remark}[theorem]{Remark}

% Common HoTT macros
\newcommand{\Type}{\mathsf{Type}}
\newcommand{\U}{\mathcal{U}}
\newcommand{\N}{\mathbb{N}}
\newcommand{\R}{\mathbb{R}}
\newcommand{\Z}{\mathbb{Z}}
\newcommand{\Q}{\mathbb{Q}}
\newcommand{\Id}{\mathsf{Id}}
\newcommand{\refl}{\mathsf{refl}}
\newcommand{\zero}{\mathsf{zero}}
\newcommand{\suc}{\mathsf{succ}}
\newcommand{\transport}{\mathsf{transport}}
\newcommand{\ua}{\mathsf{ua}}
\newcommand{\idtoeqv}{\mathsf{idtoeqv}}
% Type equivalence (\simeq) is kept as its own symbol; we provide a
% mnemonic alias \eqv that resolves to \simeq. The standard \equiv (triple
% bar) is reserved for judgmental / definitional equality and is NOT
% redefined.
\newcommand{\eqv}{\simeq}
\newcommand{\fst}{\mathsf{fst}}
\newcommand{\snd}{\mathsf{snd}}
\newcommand{\inl}{\mathsf{inl}}
\newcommand{\inr}{\mathsf{inr}}
\newcommand{\Path}{\mathsf{Path}}
% \loop is a TeX primitive; we use \looppath as a safe replacement.
\newcommand{\looppath}{\mathsf{loop}}
\newcommand{\base}{\mathsf{base}}
\newcommand{\Sone}{S^1}
\newcommand{\op}{\mathrm{op}}
\newcommand{\isContr}{\mathsf{isContr}}
\newcommand{\isProp}{\mathsf{isProp}}
\newcommand{\isSet}{\mathsf{isSet}}

% GrokRxiv sidebar configuration
\definecolor{grokgray}{RGB}{110,110,110}
\AddEverypageHook{%
  \ifnum\value{page}=1
    \begin{tikzpicture}[remember picture, overlay]
      \node[
        rotate=90,
        anchor=south,
        font=\Large\sffamily\bfseries\color{grokgray},
        inner sep=0pt
      ] at ([xshift=38pt, yshift=0.52\paperheight]current page.south west)
      {GrokRxiv:2026.05.hott-perspective\quad
       [\,math.LO\,]\quad
       03 May 2026};
    \end{tikzpicture}
  \fi
}

\title{The HoTT Perspective\\
\large Numbers as Inductive Types up to Path Equivalence}
\author{YonedaAI Research \\
\textit{Univalent Correspondence Project} \\
\textit{Magneton Foundational Studies}}
\date{3 May 2026}

\begin{document}
\maketitle

\begin{abstract}
We give a self-contained development of \emph{Homotopy Type Theory} (HoTT) as
the synthetic language in which natural numbers, real numbers, and the
transcendentals $\pi$ and $e$ admit an encoding-free description. The natural
numbers are the inductive type $\N$ generated by $\zero : \N$ and
$\suc : \N \to \N$; the integer $58$ is the canonical term
$\suc^{58}\,\zero$. Identity types $a =_A b$ are the central novelty: rather
than mere propositions, they carry the homotopical content of \emph{path
spaces}. We develop path induction, transport, and coherence, then state
Voevodsky's univalence axiom $(A = B) \eqv (A \eqv B)$ and prove its
fundamental consequence: the type $\sum_{X : \U}(X \eqv \N)$ of
``structures equivalent to $\N$'' is contractible, so every natural number
object is canonically equal as a type. We treat higher inductive types,
illustrated by the circle $\Sone{}$ with $\pi_1(\Sone{}) = \Z$, and Booij's
higher inductive--inductive construction of the Cauchy reals as an h-set,
inside which $\pi$ and $e$ are defined as the centres of contractible types
of solutions to characteristic structural conditions, formalising
``unique existence.'' Cubical type theory
(Cohen--Coquand--Huber--M\"ortberg) makes univalence computational, while
two-level, modal, and directed HoTT (Riehl--Shulman) extend the framework to
synthetic $(\infty, 1)$-category theory. We close by identifying HoTT as the
internal language of $(\infty,1)$-toposes, making the categorical and
structural views of arithmetic literally co-extensive. A companion Haskell
implementation models paths, $\refl$, $\transport$, $\mathsf{sym}$ and
$\mathsf{trans}$ at the type level via GADTs.
\end{abstract}

\tableofcontents

\section{Introduction}
\label{sec:intro}

In the previous papers of this series we asked ``What is $58$?'' and answered
it on four levels: as a numeral (Paper~01), as a von Neumann ordinal
(Paper~02), as the value of a unique morphism out of a Natural Number Object
(Paper~03), and as the representable presheaf $\mathrm{Hom}(-, 58)$
(Paper~04). Each refinement traded ontological substance for relational
structure. The question that survives Paper~04 is whether the categorical
viewpoint is \emph{itself} an encoding -- whether ``58 as the representable
$\mathrm{Hom}(-,58)$'' replaces one substantialist commitment with another.

Homotopy type theory answers ``no'' constructively. In HoTT every type is
already identified up to equivalence \emph{by axiom}, so the question
``which set, presheaf, or initial object is the real $\N$?'' is dissolved
into the slogan
\begin{quote}
\centering
\emph{equivalent things are equal.}
\end{quote}
Two natural number objects, two presentations of the reals, two $K(G,1)$'s
for the same group are not merely isomorphic up to a chosen comparison map
-- they are paths in the universe.

This paper develops the perspective in detail. We restrict ourselves to the
\emph{Book HoTT} of~\cite{HoTTBook} except where noted, and present the
material in standard Martin-L\"of style with one universe $\U$ for clarity;
the ramified hierarchy $\U_0 : \U_1 : \U_2 : \cdots$ is needed for
foundational coherence but plays no role in the arithmetic content.

\subsection{Roadmap}

Section~\ref{sec:framework} fixes notation for dependent type theory.
Section~\ref{sec:nat} defines $\N$ as an inductive type and recovers
primitive recursion. Section~\ref{sec:identity} introduces identity types,
$J$, transport, and the path groupoid. Section~\ref{sec:univalence} states
the univalence axiom and proves the contractibility of the space of
NNO-equivalents. Section~\ref{sec:hits} treats higher inductive types: the
circle, $\pi_1(\Sone{}) = \Z$, propositional truncation, the Cauchy reals.
Section~\ref{sec:cubical} sketches cubical type theory.
Section~\ref{sec:variants} surveys two-level, modal, and directed HoTT\@.
Section~\ref{sec:transcendentals} places $\pi$ and $e$ inside the Cauchy
reals. Section~\ref{sec:bridge} identifies HoTT as the internal language of
$(\infty,1)$-toposes. Section~\ref{sec:results}--\ref{sec:conclusion} record
the principal results and discuss limitations.

\subsection{Notational conventions}
\label{sec:conventions}

We write $a : A$ for ``$a$ is a term of type $A$,'' $f : A \to B$ for the
type of (non-dependent) functions, $\prod_{x:A} B(x)$ for the dependent
product (``$\Pi$-type''), and $\sum_{x:A} B(x)$ for the dependent sum
(``$\Sigma$-type''). The identity type of $a, b : A$ is written
$a =_A b$ or $\Id_A(a, b)$. The constant path is $\refl_a : a =_A a$.
Equivalence is $A \eqv B$. We use $\U$ for ``the'' universe. We deliberately
distinguish three kinds of ``equality'':
\begin{itemize}
\item the judgmental / definitional equality (a meta-level notion), written
  $a \equiv b$, used in computation rules;
\item the propositional / identity-type equality, written $a = b$, an
  inhabitant of the type $\Id_A(a, b)$;
\item type equivalence, written $A \eqv B$, an inhabitant of
  $\sum_{f : A \to B} \mathsf{isEquiv}(f)$.
\end{itemize}
Univalence asserts that the canonical map relating the second and third (at
the universe level) is itself an equivalence.

\section{Mathematical Framework}
\label{sec:framework}

We assume Martin-L\"of's intensional dependent type theory MLTT$_\omega$ as
the substrate~\cite{MartinLof84}. The pertinent rules are summarized below
informally; for the formal calculus see \cite[Appendix~A]{HoTTBook}.

\subsection{Dependent products and sums}

For a type family $B : A \to \U$, the $\Pi$-type $\prod_{x : A} B(x)$ has
introduction by $\lambda$-abstraction and elimination by application.
The $\Sigma$-type $\sum_{x : A} B(x)$ has introduction $(a, b)$ for
$a : A, b : B(a)$ and elimination by projections $\fst$ and $\snd$ together
with a dependent eliminator $\mathsf{ind}_\Sigma$.

\begin{definition}[Function type]
\label{def:fn}
$A \to B$ abbreviates $\prod_{x : A} B$ when $B$ does not depend on $x$.
\end{definition}

\begin{definition}[Cartesian product]
\label{def:prod}
$A \times B$ abbreviates $\sum_{x : A} B$ when $B$ does not depend on $x$.
\end{definition}

\subsection{The unit, empty, and sum types}

The unit type $\mathbf{1}$ has a single constructor $\star : \mathbf{1}$ and
the eliminator ``everything depending on $\mathbf{1}$ is determined by its
value at $\star$.'' The empty type $\mathbf{0}$ has no constructors and the
eliminator $\mathsf{ind}_{\mathbf{0}} : \prod_{C : \mathbf{0} \to \U}
\prod_{z : \mathbf{0}} C(z)$ (ex falso). The sum type $A + B$ has
$\inl : A \to A + B$ and $\inr : B \to A + B$.

\subsection{Universes and Russell vs.\ Tarski}
\label{sec:universes}

We adopt \emph{Russell-style} universes: $\U_0 : \U_1 : \cdots$ with
$\U_i \subset \U_{i+1}$ and a term of $\U$ is itself a type. All ``$\U$''s
in this paper denote a generic universe; statements are true at every level.

\section{Natural Numbers as an Inductive Type}
\label{sec:nat}

\subsection{The inductive definition}

\begin{definition}[Natural numbers]
\label{def:nat}
The type $\N : \U$ is generated by the formation rule $\N : \U$ and the two
constructors
\begin{align*}
\zero &: \N, \\
\suc &: \N \to \N.
\end{align*}
The dependent eliminator (\emph{induction principle}) is
\[
\mathsf{ind}_\N : \prod_{P : \N \to \U}\;
P(\zero) \to
\Bigl(\prod_{n : \N} P(n) \to P(\suc\,n)\Bigr) \to
\prod_{n : \N} P(n).
\]
The computation rules are
$\mathsf{ind}_\N\,P\,c_0\,c_s\,\zero \equiv c_0$ and
$\mathsf{ind}_\N\,P\,c_0\,c_s\,(\suc\,n) \equiv
c_s\,n\,(\mathsf{ind}_\N\,P\,c_0\,c_s\,n)$.
\end{definition}

\begin{definition}[Numerals]
\label{def:numerals}
Define $\overline{n} : \N$ recursively by $\overline{0} := \zero$ and
$\overline{n+1} := \suc\,\overline{n}$. In particular,
\[
\overline{58} = \suc(\suc(\suc(\cdots(\suc\,\zero)\cdots)))
\quad\text{(58 applications of \(\suc\)).}
\]
\end{definition}

\subsection{Recursion and arithmetic}

The non-dependent recursor
\[
\mathsf{rec}_\N : \prod_{C : \U} C \to (\N \to C \to C) \to \N \to C
\]
is the special case where $P(n) := C$ for all $n$.

\begin{example}[Addition]
\label{ex:add}
Define $+ : \N \to \N \to \N$ by
$m + n := \mathsf{rec}_\N\,\N\,n\,(\lambda k.\,\suc)\,m$, equivalently by
$\zero + n \equiv n$ and $(\suc\,m) + n \equiv \suc(m + n)$. Multiplication
and exponentiation are defined analogously.
\end{example}

\begin{proposition}[NNO universal property, internal]
\label{prop:nno}
For any type $C : \U$, point $c_0 : C$, and endo-map $f : C \to C$, there is
a unique $h : \N \to C$ with $h(\zero) = c_0$ and
$h(\suc\,n) = f(h\,n)$ for all $n : \N$.
\end{proposition}

\begin{proof}
Existence: $h := \mathsf{rec}_\N\,C\,c_0\,(\lambda n.\,f)$. By the
recursor's computation rules we have the \emph{judgmental} equalities
$h(\zero) \equiv c_0$ and $h(\suc\,n) \equiv f(h\,n)$.

Uniqueness: suppose $h' : \N \to C$ is given together with paths
$p_0 : h'(\zero) = c_0$ and
$p_s : \prod_{n : \N} h'(\suc\,n) = f(h'\,n)$
witnessing that $h'$ also satisfies the equations (note these are
\emph{propositional} equalities -- judgmental equality cannot be assumed
for a generic $h'$). Define the predicate
$P(n) := (h\,n =_C h'\,n)$.

\emph{Base case.} Since $h(\zero) \equiv c_0$ judgmentally, the
inverse path $p_0^{-1} : c_0 = h'(\zero)$ inhabits $P(\zero)$.

\emph{Inductive step.} Assume $q : h\,n = h'\,n$. The recursor's
computation rule gives $h(\suc\,n) \equiv f(h\,n)$ judgmentally, and
applying $\mathsf{ap}_f : (a = b) \to (f\,a = f\,b)$ to $q$ yields
$\mathsf{ap}_f(q) : f(h\,n) = f(h'\,n)$. Composing with $p_s(n)^{-1}$
gives $f(h\,n) = h'(\suc\,n)$, hence $h(\suc\,n) = h'(\suc\,n)$.
Apply $\mathsf{ind}_\N$ to conclude $\prod_{n : \N} P(n)$. Function
extensionality (\Cref{thm:funext}) then gives $h = h'$.
\end{proof}

\subsection{$\N$ in HoTT vs.\ $\N$ in set theory}
\label{sec:nat-vs-set}

\begin{remark}
\label{rem:benacerraf}
Proposition~\ref{prop:nno} renders Benacerraf's identification problem moot
\emph{within} the type theory: $\N$ is determined up to canonical bijection
by its formation and elimination rules. Combined with univalence
(\Cref{sec:univalence}) it is determined up to canonical \emph{equality} as a
term of $\U$.
\end{remark}

\section{Identity Types and Path Induction}
\label{sec:identity}

\subsection{The identity type}

\begin{definition}[Identity type]
\label{def:id}
For each type $A : \U$ and elements $a, b : A$ there is a type
$a =_A b$, called the \emph{identity type} or \emph{path space}. It is the
inductive family generated by the single constructor
$\refl_a : a =_A a$ for each $a : A$.
\end{definition}

The eliminator -- often called \emph{path induction} or the \emph{$J$-rule}
-- says that to prove a property of all paths it suffices to prove it of the
constant paths.

\begin{proposition}[Path induction, $J$]
\label{prop:J}
For every type $A : \U$ and family
$C : \prod_{x, y : A} (x =_A y) \to \U$, given a section
$d : \prod_{x : A} C(x, x, \refl_x)$, there is a section
\[
J(C, d) : \prod_{x, y : A} \prod_{p : x =_A y} C(x, y, p),
\]
satisfying the computation rule $J(C, d, x, x, \refl_x) \equiv d(x)$.
\end{proposition}

\begin{remark}
\label{rem:J-vs-K}
$J$ is \emph{not} the principle ``every path is $\refl$.'' That stronger
statement is Streicher's $K$-axiom and is incompatible with univalence: it
forces every type to be a set. HoTT keeps $J$ but rejects $K$.
\end{remark}

\subsection{Symmetry, transitivity, and the path groupoid}

\begin{lemma}[Symmetry]
\label{lem:sym}
There is a function
$\mathsf{sym} : \prod_{A : \U}\prod_{x, y : A} (x =_A y) \to (y =_A x)$.
\end{lemma}

\begin{proof}
Take $C(x, y, p) := (y =_A x)$ and $d(x) := \refl_x$, then apply $J$.
\end{proof}

\begin{lemma}[Transitivity]
\label{lem:trans}
There is a function $\mathsf{trans} : \prod_{A : \U}\prod_{x, y, z : A}
(x =_A y) \to (y =_A z) \to (x =_A z)$.
\end{lemma}

\begin{proof}
Fix $A, x, y, z$ and $p : x = y$. By $J$ on $p$ it suffices to assume
$y \equiv x$ and $p \equiv \refl_x$, in which case
$\mathsf{trans}(\refl_x, q) := q$ is a section.
\end{proof}

We write $p^{-1}$ for $\mathsf{sym}(p)$ and $p \cdot q$ for
$\mathsf{trans}(p, q)$.

\begin{theorem}[Path groupoid]
\label{thm:groupoid}
For any type $A$ and $a, b, c, d : A$, $p : a = b$, $q : b = c$, $r : c = d$:
\begin{enumerate}
\item $\refl_a \cdot p = p$ and $p \cdot \refl_b = p$;
\item $p \cdot p^{-1} = \refl_a$ and $p^{-1} \cdot p = \refl_b$;
\item $(p \cdot q) \cdot r = p \cdot (q \cdot r)$.
\end{enumerate}
Moreover the higher coherences (the pentagon, the Eckmann--Hilton
argument, etc.) all hold up to higher paths, making $A$ into an
$\infty$-groupoid~\cite{LumsdaineGroupoid, vandenBergGarner}.
\end{theorem}

\begin{proof}[Proof sketch]
Each clause is proved by $J$ on the relevant paths. For the higher coherences
one iterates $J$ at successively higher levels; see
\cite[\S2.1]{HoTTBook}.
\end{proof}

\subsection{Transport}

\begin{definition}[Transport]
\label{def:transport}
For $P : A \to \U$, $a, b : A$, and $p : a = b$, the function
$\transport^P(p) : P(a) \to P(b)$ is defined by $J$:
$\transport^P(\refl_a)(u) := u$.
\end{definition}

\begin{proposition}[Transport along concatenation]
\label{prop:trans-concat}
\[
\transport^P(p \cdot q) = \transport^P(q) \circ \transport^P(p).
\]
\end{proposition}

\begin{proof}
$J$ on either $p$ or $q$.
\end{proof}

\subsection{Homotopy levels}

\begin{definition}[Contractibility, propositions, sets]
\label{def:hlevel}
A type $A : \U$ is
\begin{itemize}
\item \emph{contractible}, $\isContr(A) := \sum_{c : A} \prod_{x : A} (x = c)$;
\item a \emph{mere proposition}, $\isProp(A) := \prod_{x, y : A} (x = y)$;
\item a \emph{set} (h-set), $\isSet(A) := \prod_{x, y : A} \isProp(x = y)$.
\end{itemize}
The hierarchy continues: $A$ is an $n$-type iff every identity type of $A$
is an $(n-1)$-type.
\end{definition}

\begin{theorem}[Hedberg]
\label{thm:hedberg}
If $A$ has decidable equality $\prod_{x, y : A} (x = y) + \neg(x = y)$, then
$A$ is a set.
\end{theorem}

\begin{corollary}
\label{cor:N-set}
$\N$ is a set.
\end{corollary}

\begin{proof}
Equality on $\N$ is decidable by induction on both arguments and Peano's
fourth axiom $\suc\,m = \suc\,n \Rightarrow m = n$. Apply Hedberg.
\end{proof}

\section{The Univalence Axiom}
\label{sec:univalence}

\subsection{Equivalence of types}

\begin{definition}[Quasi-inverse and equivalence]
\label{def:equiv}
A function $f : A \to B$ has a \emph{quasi-inverse} if there exist
$g : B \to A$ and homotopies $\eta : g \circ f \sim \mathrm{id}_A$ and
$\varepsilon : f \circ g \sim \mathrm{id}_B$. The proposition
$\mathsf{isEquiv}(f)$ is the half-adjoint version
\cite[\S4.2]{HoTTBook}, ensuring $\mathsf{isEquiv}(f)$ is itself a mere
proposition. We write $A \eqv B := \sum_{f : A \to B} \mathsf{isEquiv}(f)$.
\end{definition}

\begin{proposition}[idtoeqv]
\label{prop:idtoeqv}
There is a canonical map
$\idtoeqv : (A =_\U B) \to (A \eqv B)$ taking $\refl_A$ to the identity
equivalence.
\end{proposition}

\begin{proof}
$J$ on the path $A = B$.
\end{proof}

\subsection{The axiom}

\begin{theorem}[Voevodsky's univalence axiom]
\label{thm:univalence}
\textup{(\cite{Voevodsky}; \cite[Axiom 2.10.3]{HoTTBook}.)}
For all types $A, B : \U$, the canonical map $\idtoeqv$ is itself an
equivalence:
\[
\ua : (A \eqv B) \to (A =_\U B), \qquad
\idtoeqv \circ \ua = \mathrm{id}_{A \eqv B}, \qquad
\ua \circ \idtoeqv = \mathrm{id}_{A =_\U B}.
\]
\end{theorem}

In Book HoTT this is taken as an axiom; in cubical type theory
(\Cref{sec:cubical}) it is a theorem.

\subsection{Consequences for arithmetic}

\begin{definition}[Type of NNOs]
\label{def:nno-type}
Define
\[
\mathrm{NNO} := \sum_{X : \U}\;\sum_{x_0 : X}\;\sum_{s : X \to X}\;
\mathsf{isInitial}(X, x_0, s),
\]
where $\mathsf{isInitial}(X, x_0, s)$ is the universal property of
\Cref{prop:nno} expressed as a $\Sigma$-type.
\end{definition}

\begin{theorem}[Contractibility of NNO]
\label{thm:nno-contractible}
$\mathrm{NNO}$ is contractible.
\end{theorem}

\begin{proof}
We exhibit $(\N, \zero, \suc, \mathsf{init}_\N)$ as the centre and supply
a proof that every other quadruple is equal to it.

\emph{Step 1. Equivalence at the carrier.}
By \Cref{prop:nno} the standard structure $(\N, \zero, \suc)$ is initial.
For any other initial $(X, x_0, s)$, the universal property applied in
both directions yields mutually inverse maps $f : \N \to X$ and
$g : X \to \N$ with $g \circ f = \mathrm{id}_\N$ and
$f \circ g = \mathrm{id}_X$ (each by uniqueness of the universal
recursor). Half-adjointness packages this as an equivalence
$e : \N \eqv X$. Univalence (\Cref{thm:univalence}) promotes this to a
path $p := \ua(e) : \N =_\U X$.

\emph{Step 2. Paths in $\Sigma$-types.}
Recall that a path in $\sum_{a : A} B(a)$ between $(a, b)$ and
$(a', b')$ is, up to canonical isomorphism, a pair $(p, q)$ where
$p : a =_A a'$ and $q : \transport^B(p)(b) =_{B(a')} b'$
\cite[Theorem~2.7.2]{HoTTBook}. Iterating: a path between two iterated
$\Sigma$-tuples
$(a_1, a_2, \ldots, a_n)$ and $(a_1', a_2', \ldots, a_n')$ is a tuple of
paths $(p_1, p_2, \ldots, p_n)$ where $p_k$ lives \emph{over} the
preceding $p_1, \ldots, p_{k-1}$.

\emph{Step 3. Constructing the path componentwise.}
We must produce paths $(p, p_0, p_s, p_*)$ where:
\begin{enumerate}
\item $p : \N =_\U X$ is the carrier path from Step~1.
\item $p_0 : \transport^{(X \mapsto X)}(p)(\zero) =_X x_0$. Because
$\ua(e)$ transports along the equivalence $e$, the left-hand side
reduces to $e(\zero)$, and the universal property of $\N$ forces
$e(\zero) = x_0$ (the unique map $\N \to X$ sends $\zero$ to $x_0$),
giving the required path.
\item $p_s : \transport^{(X \mapsto X \to X)}(p)(\suc) = s$.
The type family $X \mapsto X \to X$ has transport given by conjugation
by $e$, so the left-hand side reduces to $e \circ \suc \circ e^{-1}$.
Naturality of the universal recursor yields $e \circ \suc = s \circ e$,
whence $e \circ \suc \circ e^{-1} = s$, supplying $p_s$.
\item $p_*$: the initiality witness $\mathsf{isInitial}$ is a mere
proposition (uniqueness of the universal recursor is a Pi-type into a
contractible identity type), so any two initiality proofs are
identified by a unique path.
\end{enumerate}

\emph{Step 4. Conclusion.}
Bundling $(p, p_0, p_s, p_*)$ via the iterated $\Sigma$-path
constructor gives the required equality
$(\N, \zero, \suc, \mathsf{init}_\N) =_{\mathrm{NNO}} (X, x_0, s, \cdot)$.
Since $(X, x_0, s, \cdot)$ was arbitrary, $\mathrm{NNO}$ is contractible.
The argument is the same shape as the contractibility of the type of
``small structures-of-shape-$T$ on a fixed base'' once $T$ is replaced
by the signature ``one constant and one unary operation''; see
\cite[\S5.4 and \S2.7]{HoTTBook}.
\end{proof}

\begin{corollary}
\label{cor:58-canonical}
For any NNO $(X, x_0, s)$ the term $s^{58}\,x_0$ is, after transport along
the canonical equality $\N = X$, equal to $\overline{58} : \N$.
\end{corollary}

\Cref{cor:58-canonical} is the formal expression of the slogan ``58 is
canonical across all NNOs.''

\subsection{Function extensionality}

\begin{theorem}[Function extensionality from univalence]
\label{thm:funext}
If $f, g : \prod_{x : A} B(x)$ and there is a homotopy
$h : \prod_{x : A} (f\,x = g\,x)$, then $f = g$.
\end{theorem}

\begin{proof}[Proof sketch]
Voevodsky proved this by representing $\prod_{x : A} B(x)$ via a clever
$\Sigma$-type and applying univalence. See \cite[\S4.9]{HoTTBook}.
\end{proof}

\section{Higher Inductive Types}
\label{sec:hits}

A higher inductive type (HIT) extends inductive types with constructors for
\emph{paths} as well as points.

\subsection{The interval and propositional truncation}

\begin{definition}[Interval]
The type $I$ has constructors $0_I, 1_I : I$ and a path
$\mathsf{seg} : 0_I = 1_I$.
\end{definition}

\begin{definition}[Propositional truncation]
\label{def:trunc}
$\|A\|$ has $|a| : \|A\|$ for $a : A$ and a path constructor
$\prod_{x, y : \|A\|}(x = y)$. Its eliminator allows mapping $\|A\| \to P$
when $P$ is a mere proposition.
\end{definition}

\subsection{The circle}

\begin{definition}[Circle]
\label{def:circle}
$\Sone$ has constructors $\base : \Sone$ and a path
$\looppath : \base =_{\Sone} \base$. The eliminator: given $P : \Sone \to \U$,
$b : P(\base)$, and $\ell : \transport^P(\looppath)(b) = b$, there is
$\mathsf{ind}_{\Sone}(P, b, \ell) : \prod_{x : \Sone} P(x)$, computing as
$\mathsf{ind}_{\Sone}(P, b, \ell)(\base) \equiv b$ and a corresponding rule
on $\looppath$.
\end{definition}

\begin{theorem}[Loop space of \texorpdfstring{$\Sone{}$}{S\^{}1}]
\label{thm:pi1}
$\pi_1(\Sone) := \|(\base = \base)\|_0 \eqv \Z$.
Indeed $(\base = \base) \eqv \Z$ as a type, before truncation.
\end{theorem}

\begin{proof}[Proof sketch]
Define the universal cover
$\mathsf{code} : \Sone \to \U$ by $\mathsf{code}(\base) := \Z$ and on
$\looppath$ by univalence and the successor equivalence
$\mathsf{succ} : \Z \eqv \Z$. Then
$(\base = -) \sim \mathsf{code}$ by the encode--decode method
(Licata--Shulman~\cite{LicataShulman}); evaluating at $\base$ gives
$(\base = \base) \eqv \Z$.
\end{proof}

\Cref{thm:pi1} is the first non-trivial theorem of \emph{synthetic homotopy
theory}: a result of algebraic topology proved entirely inside type theory,
without recourse to topological spaces, simplicial sets, or model
categories.

\subsection{Spheres, suspensions, and pushouts}

The $n$-sphere $S^n$ is defined recursively as $\Sigma S^{n-1}$ where
$\Sigma$ is the unreduced suspension HIT\@. The pushout of
$f : C \to A$ and $g : C \to B$ has constructors for $A$, $B$, and a path
identifying $f(c)$ and $g(c)$ for each $c : C$.

\begin{theorem}[Hopf fibration, internal]
\label{thm:hopf}
There is a fibration $h : S^3 \to S^2$ with fibers $\Sone{}$, defined
internally as a type family $H : S^2 \to \U$ via the eliminator of the
HIT $S^2$ (with total space $\sum_{x : S^2} H(x) \eqv S^3$), and
characterized by the long exact sequence of homotopy groups; in
particular $\pi_3(S^2) \eqv \Z$ via the Hopf invariant
\cite[\S8.5]{HoTTBook}.
\end{theorem}

\subsection{Eilenberg--MacLane spaces}

For each abelian group $G$ and $n \geq 1$ there is a HIT $K(G, n)$ with
$\pi_n = G$ and $\pi_k = 0$ for $k \neq n$
\cite{LicataFinster}. Synthetic cohomology
$H^n(X; G) := \|X \to K(G, n)\|_0$ recovers ordinary singular cohomology
when $X$ is a CW-complex.

\subsection{Real numbers as a higher inductive--inductive type}
\label{sec:reals}

\begin{theorem}[Booij; HoTT Book \S11.3]
\label{thm:reals}
\textup{(\cite{Booij}, after~\cite[\S11.3]{HoTTBook}.)}
There is a higher inductive--inductive type $\R_c$ presenting the
\emph{Cauchy reals}: an h-set with constructors for rational injections
$\mathsf{rat} : \Q \to \R_c$, a Cauchy-completion limit operation, and a
path constructor identifying terms whose distance witnesses are below every
positive rational. $\R_c$ has the universal property of the Cauchy
completion of $\Q$ in the category of metric spaces internal to HoTT, and
satisfies the axioms of an Archimedean ordered field.
\end{theorem}

The construction is \emph{higher} because of the path constructor and
\emph{inductive--inductive} because the path constructor mentions a
``closeness'' relation defined simultaneously with the carrier type.

\subsection{Worked example: the canonical 58 across NNOs}
\label{sec:worked-example}

To make \Cref{thm:nno-contractible} concrete, we walk through how an
arbitrary NNO yields the same ``58'' as $\N$.

\begin{example}[58 in two NNOs]
\label{ex:58-two-nnos}
Take two structures.
\begin{itemize}
\item The standard one: $(\N, \zero, \suc)$.
\item A bit-string encoding: $X := \mathrm{List}(\mathbf{1})$, with
$x_0 := \mathsf{nil}$ and $s := \mathsf{cons}(\star, -)$.
\end{itemize}
Both are easily checked to satisfy the universal property of
\Cref{prop:nno}; for the bit-string version, the unique recursor sends
a list of length $n$ to the value $f^n(c_0)$ in any target $(C, c_0, f)$.

By \Cref{thm:nno-contractible}, the data $(\N, \zero, \suc)$ and
$(X, x_0, s)$ are equal as inhabitants of $\mathrm{NNO}$. Concretely:
\begin{enumerate}
\item the equivalence $e : \N \eqv X$ sends $\overline{n}$ to the list
of length $n$;
\item univalence gives $\ua(e) : \N =_\U X$;
\item transporting $\overline{58} : \N$ along $\ua(e)$ yields the
58-element list of $\star$'s in $X$;
\item this list is, by definition of $X$, the value of
$s^{58}\,x_0$.
\end{enumerate}
Thus the two ``58''s are propositionally equal, and \emph{this very
equality} is the content of \Cref{cor:58-canonical}.
\end{example}

\begin{remark}[The Yoneda perspective recovered]
\label{rem:yoneda-as-univalence}
Paper~04 characterises every object $A$ by its representable functor
$\mathrm{Hom}(-, A)$ and shows $A \simeq B \Leftrightarrow
\mathrm{Hom}(-, A) \simeq \mathrm{Hom}(-, B)$ (Yoneda). Univalence
upgrades the right-to-left direction: \emph{equivalent representables
yield equal objects in $\U$}, which is the type-theoretic statement of
the structuralist slogan ``a thing is what it relates to.''
\end{remark}

\subsection{Worked example: the loop concatenation as integer addition}
\label{sec:s1-example}

For \Cref{thm:pi1}, let us trace what happens to specific loops.

\begin{example}[Three full turns then two backward]
\label{ex:loops}
Inside HoTT we represent the loop ``three times around'' as the path
$\looppath \cdot \looppath \cdot \looppath$ at $\base$ in $\Sone$, i.e.\ a term
of $\base = \base$. By the encode--decode equivalence
$(\base = \base) \eqv \Z$, this term encodes to the integer $3$.

Composing with $\looppath^{-1} \cdot \looppath^{-1}$ gives, by
\Cref{thm:groupoid}, the path
$\looppath \cdot \looppath \cdot \looppath \cdot \looppath^{-1} \cdot \looppath^{-1}$,
which the encode-decode equivalence sends to $3 + (-1) + (-1) = 1$.

Decoding back yields a single $\looppath$, propositionally equal to the
original concatenation. The point is not that we computed $3 - 2 = 1$
(routine!), but that this very arithmetic of $\Z$ \emph{is} the
fundamental group law on $\pi_1(\Sone)$ once we identify the two via
the equivalence.
\end{example}

\section{Cubical Type Theory}
\label{sec:cubical}

The original presentation of HoTT takes univalence as an axiom, breaking
canonicity: a closed term of $\N$ may be propositionally but not
definitionally equal to a numeral. Cubical type theory restores
canonicity.

\subsection{Intervals and faces}

\begin{definition}[Cube category]
\label{def:cube}
Cohen--Coquand--Huber--M\"ortberg~\cite{CCHM} introduce a primitive
interval object $\mathbb{I}$ generated by $0, 1 : \mathbb{I}$, the
operations $\sqcup, \sqcap, \neg$ making $\mathbb{I}$ into a de Morgan
algebra, and the algebra of \emph{face formulas} $\varphi$ over the
generating equations $i = 0$, $i = 1$.
\end{definition}

A path $p : a =_A b$ is then a function $\mathbb{I} \to A$ with $p(0) = a$
and $p(1) = b$ judgmentally. Path concatenation, inversion, and dependent
elimination become explicit term operations on cubical sets, not derived
from $J$.

\subsection{Composition and Glue}

The crucial operations are \emph{composition} (filling open boxes) and
\emph{Glue} (gluing types along an equivalence). The Glue type
$\mathsf{Glue}\,A\,[\varphi \mapsto (T, f)]$ where $f : T \eqv A$ on the
face $\varphi$ supplies a constructive proof of univalence: every
equivalence is realized as a path in $\U$.

\begin{theorem}[Univalence in cubical type theory]
\label{thm:cubical-univalence}
Cubical type theory in the CCHM model proves $\ua$ as a defined term, and
$\idtoeqv \circ \ua$ reduces to the identity by computation, not by
postulate~\cite{CCHM, CubicalAgda}.
\end{theorem}

\subsection{Implementations}

The Cubical Agda mode~\cite{CubicalAgda} and the
\texttt{cubicaltt} prototype implement the system; the \texttt{redtt} and
\texttt{cooltt} projects extend it with cartesian intervals and modalities.
For our purposes the relevant fact is that all theorems of
\Cref{sec:univalence,sec:hits} can be \emph{computed}, not just postulated.

\section{Variants: 2LTT, Modal, and Directed HoTT}
\label{sec:variants}

\subsection{Two-level type theory}

Two-level type theory (Annenkov--Capriotti--Kraus--Sattler~\cite{2LTT})
introduces a strict equality alongside the homotopical one. The strict
fragment is used to formalize semantics (e.g.\ inverse diagrams of types
representing semisimplicial types) that the homotopical fragment alone
cannot encode without higher coherence. Strict equality is incompatible
with univalence, so 2LTT carefully separates the two universes.

\subsection{Modal HoTT}

Modal type theory adds operators $\square, \diamondsuit$ corresponding to
adjunctions in the categorical semantics. Examples include
Schreiber--Shulman cohesive HoTT~\cite{SchreiberShulman}, where the
cohesion modalities recover the discrete--codiscrete--shape adjoint string
$\mathrm{disc} \dashv \Pi \dashv \mathrm{codisc}$ for cohesive
$(\infty,1)$-toposes, and the spatial type theory of Riley
\cite{Riley}.

\subsection{Directed HoTT}

Directed type theory replaces invertible paths with possibly non-invertible
\emph{morphisms}, so that types correspond to $(\infty, n)$-categories
rather than $\infty$-groupoids. Riehl--Shulman~\cite{RiehlShulman} develop
\emph{simplicial type theory} with a primitive directed interval $\Delta^1$
and prove a synthetic Yoneda lemma; directed univalence
\cite{CisinskiNguyen} is an active research frontier.

\section{Transcendentals in HoTT}
\label{sec:transcendentals}

We work inside $\R_c$ from \Cref{sec:reals}.

\subsection{\texorpdfstring{$\pi$}{pi} as the period of \texorpdfstring{$\exp$}{exp} on \texorpdfstring{$i\R$}{iR}}

\begin{definition}[Real exponential]
The map $\exp : \R_c \to \R_c$ is defined as the unique solution to
$y' = y, y(0) = 1$ in $\R_c$ -- equivalently, by the power series
$\sum_{n=0}^{\infty} x^n / n!$, which converges in $\R_c$ by completeness.
\end{definition}

\begin{definition}[Pi]
\label{def:pi}
Define $\sin : \R_c \to \R_c$ by the power series
$\sum_{k=0}^\infty (-1)^k x^{2k+1}/(2k+1)!$. The unique-existence
proposition
\[
\sum_{p : \R_c}\;\bigl(\sin(p) = 0\bigr) \times \bigl(p > 0\bigr) \times
\bigl(\prod_{x : \R_c}\;0 < x < p \to \sin(x) > 0\bigr)
\]
is contractible. We define $\pi$ to be its centre.
\end{definition}

\begin{theorem}
\label{thm:pi-trans}
Lindemann's theorem ports to HoTT: $\pi$ is transcendental over $\Q$, in
the sense that the type
\[
\sum_{p : \mathrm{Poly}(\Q)}\;\bigl(p \neq 0\bigr) \times
\bigl(p(\pi) = 0\bigr)
\]
is empty (equivalent to $\mathbf{0}$).
\end{theorem}

\begin{proof}[Proof sketch]
Lindemann's classical proof relies only on convergent power series,
elementary symmetric functions, and bounding arguments -- all of which
are available constructively in $\R_c$. The port is therefore believed
to be possible based on the constructive character of the classical
argument; however, to the authors' knowledge no end-to-end formalization
in HoTT is yet complete. The presently available scaffolding includes:
\begin{itemize}
\item Cauchy completeness of $\R_c$, basic limits, the Archimedean
property, and elementary calculus on $\R_c$ (Booij's thesis
\cite{Booij}, expanding \cite[\S11.3]{HoTTBook}).
\item Convergence of power series and analytic functions in HoTT, plus
algebraic manipulations of polynomial rings and symmetric functions
(scaffolded in the agda-unimath project~\cite{AgdaUnimath}).
\end{itemize}
The pieces still missing for an end-to-end formalization are: (i)~the
Hermite--Lindemann auxiliary function $J(t) = \int_0^t e^{t - s}
f(s)\,ds$ and its bounds in $\R_c$; (ii)~the algebraic-integer estimates
that drive the contradiction; and (iii)~the integration with a finished
formalization of the Cauchy integral over $\R_c$. We identify this as
open work in \Cref{sec:conclusion}.
\end{proof}

\subsection{$e$ as the unique fixed point of $D$}

\begin{definition}[Euler's number]
\label{def:e}
Let $\exp : \R_c \to \R_c$ be as above. Then $e := \exp(1)$. Equivalently,
$e$ is the centre of the contractible type
$\sum_{x : \R_c}\;\sum_{f : \R_c \to \R_c}\;
(f' = f) \times (f(0) = 1) \times (x = f(1))$.
\end{definition}

\begin{theorem}[Hermite, ported]
\label{thm:e-trans}
$e$ is transcendental over $\Q$.
\end{theorem}

\subsection{Identification across encodings}

\begin{theorem}[Reals are unique in HoTT up to equivalence]
\label{thm:reals-unique}
Any two h-sets satisfying the universal property of the Cauchy reals (or
the Dedekind reals, modulo the law of excluded middle for the latter) are
equivalent, hence equal as types by univalence. Consequently, $\pi$ and $e$
are well-defined independently of the chosen construction.
\end{theorem}

\section{Bridge: HoTT as the Internal Language of \texorpdfstring{$(\infty,1)$}{(infinity,1)}-Toposes}
\label{sec:bridge}

\subsection{Categorical semantics}

\begin{theorem}[Awodey--Warren]
\label{thm:aw}
\textup{(\cite{AwodeyWarren}.)}
Identity types of intensional Martin-L\"of type theory are interpreted in
weak factorisation systems; the model in simplicial sets exhibits the path
groupoid as a Kan complex.
\end{theorem}

Voevodsky's simplicial set model~\cite{KapulkinLumsdaine} establishes that
HoTT with univalence has a sound interpretation in the
$(\infty, 1)$-topos $\mathsf{sSet}$ of Kan complexes.

\begin{theorem}[Shulman]
\label{thm:shulman-it}
\textup{(\cite{ShulmanInftyTopos}.)}
Every $(\infty, 1)$-topos $\mathcal{E}$ admits a model of HoTT with
univalence; conversely, the syntactic category of HoTT is conjecturally
the initial $(\infty, 1)$-topos.
\end{theorem}

\subsection{Co-extensiveness with categorical and structural perspectives}

The Yoneda embedding (Paper~04) and the categorical structuralism of
Paper~06 have always been ``equivalent up to equivalence.'' Univalence
makes them \emph{equal} -- equal as objects of $\U$, equal as theories
under semantic interpretation, equal as derivations modulo a translation
that is itself an equivalence.

\begin{remark}[The univalent correspondence, formalized]
The Curry--Howard--Lambek tradition (see Wadler's exposition~\cite{Wadler2015})
identifies logic, type theory, and category theory as a triple. The
four-column correspondence
\[
\text{Logic} \;\longleftrightarrow\;
\text{Type Theory} \;\longleftrightarrow\;
\text{Category Theory} \;\longleftrightarrow\;
\text{Homotopy}
\]
is, in HoTT, the statement that the four interpretations of the same term
in $\U$ correspond to the same homotopy type, with the categorical column
interpreted in $(\infty, 1)$-toposes (Lurie's higher topos
theory~\cite{Lurie} provides the ambient setting). In particular,
``58 as a proposition'' (the inhabited type
$\mathsf{Fin}_{59} \cap \mathsf{prime}^c$, etc.), ``58 as a term''
($\overline{58} : \N$), ``58 as an object'' (the discrete object
$58 \in \mathsf{Set}$), and ``58 as a space'' (the discrete homotopy type
$\mathsf{B}\mathbf{1}^{58}$) are equal up to canonical paths.
\end{remark}

\section{Results}
\label{sec:results}

We collect the principal contributions of this paper.

\begin{theorem}[Encoding-free arithmetic]
\label{thm:main}
In Book HoTT with univalence, the following are equivalent:
\begin{enumerate}
\item the inductive type $\N$ with $\zero, \suc$;
\item the initial NNO in any model of HoTT with $\N$-types;
\item the unique pointed endo-structure $X$ with $\mathsf{isInitial}(X, x_0, s)$;
\item up to univalent equality, every other type satisfying any of (1)--(3).
\end{enumerate}
The witnesses to these equivalences are themselves canonical (paths in
$\U$), so $58 := \overline{58}$ is independent of representation.
\end{theorem}

\begin{theorem}[Synthetic computation of \texorpdfstring{$\pi_1(\Sone{})$}{pi\_1(S\^{}1)}]
\label{thm:pi1-result}
Within HoTT, $\pi_1(\Sone{}) = \Z$ admits a fully formal proof
\cite{LicataShulman} that mechanizes in Agda and Cubical Agda; the proof
uses no external topology, only the HIT $\Sone$ and univalence.
\end{theorem}

\begin{theorem}[Reals and transcendentals]
\label{thm:reals-result}
The Cauchy reals $\R_c$ form an h-set inside HoTT
(\Cref{thm:reals}). The terms $\pi$ and $e$ are unique-existence centres of
contractible types of solutions of structural conditions. Their classical
properties (irrationality, transcendence) port mechanically.
\end{theorem}

\begin{theorem}[Bridge theorem]
\label{thm:bridge-result}
HoTT is sound and complete for $(\infty, 1)$-topos semantics
\cite{ShulmanInftyTopos}. Hence the categorical (Paper~06), structural
(Paper~07), Yoneda (Paper~04), and HoTT (Paper~05) views of arithmetic are
co-extensive: a statement provable in any one is provable in any other,
modulo translation.
\end{theorem}

\section{Discussion}
\label{sec:discussion}

\subsection{What HoTT does not solve}

HoTT is a foundation, not a panacea. The following remain open:

\begin{itemize}
\item \emph{Coherence beyond truncation level}: proving general coherence
theorems for $(\infty, 1)$-categories \emph{inside} HoTT (rather than via
2LTT) is non-trivial. The semisimplicial types problem
\cite{SemisimplicialProblem} is the canonical example.
\item \emph{Decidability and computation under axioms}: Book HoTT loses
canonicity. Cubical type theory restores it but at the cost of a more
intricate calculus.
\item \emph{Set-theoretic strength}: HoTT with univalence has the
proof-theoretic strength of ZFC + ``there exist many inaccessible
cardinals'' \cite{ShulmanStrength}, comparable to standard mathematical
practice. Whether all of analytic number theory ports remains an open
research front (Paper~07 surveys this).
\end{itemize}

\subsection{Philosophical reading}

The shift from set theory to HoTT mirrors the shift from substantialism to
structuralism. ``58 is a thing'' becomes ``58 is a position in a
structure'' becomes ``58 is the equivalence class of all positions in
isomorphic structures'' becomes (in HoTT) ``58 is a term whose identity is
witnessed by paths.'' Univalence collapses the apparent distinction
between the last two.

\subsection{Pedagogical implications}

Synthetic homotopy theory provides a path-from-paths approach to algebraic
topology: $\pi_1(\Sone{}) = \Z$ in twelve lines of Agda. This is not merely a
streamlining; it indicates that homotopy is more primitive than topology
in the sense that the latter can be dispensed with.

\subsection{Computational realisation}

The Haskell companion (\Cref{sec:haskell}) and the Agda fragments
(\Cref{sec:agda-sketch}) demonstrate that path induction, transport, and
the basic equational reasoning calculus can be \emph{simulated} at the
type level in mainstream functional languages, even without full
dependent types. This is a useful pedagogical bridge, not a replacement.

\subsection{Comparison with set-theoretic and categorical foundations}

It is illuminating to put the four foundational stances side by side, as
they answer the question ``what is the natural number $58$?''.

\begin{itemize}
\item \textbf{Set theory (Paper~02).} ``$58$ is the von Neumann ordinal
$\{0, 1, \ldots, 57\}$.'' The encoding is canonical only by stipulation;
Benacerraf's identification problem records the residual arbitrariness.

\item \textbf{Categorical foundations (Paper~03 / Paper~04).}
``$58$ is the value of the unique recursion morphism on the natural-number
object, equivalently a position in the representable
$\mathrm{Hom}(-, \N)$.'' Encoding-independence holds up to canonical
isomorphism; the question ``but \emph{which} NNO?'' remains a residue,
albeit a small one.

\item \textbf{HoTT (this paper).} ``$58$ is $\overline{58} : \N$, where
$\N$ is determined up to a canonical \emph{equality} of types by
univalence.'' Encoding-independence is internal: there is no extra-
theoretic notion of ``which $\N$.''

\item \textbf{Structural / $\infty$-topos perspective (Paper~06).} ``$58$
is an invariant under all structure-preserving maps,'' which by the
internal-language theorem is the same as the HoTT statement.
\end{itemize}

The progression is one of \emph{decreasing residue}: each step removes a
notion that the previous level could not eliminate.

\subsection{When axiomatic univalence cannot be avoided}

In Book HoTT, univalence is \emph{not} a constructive principle: a closed
term of $\N$ produced via univalence may fail to reduce to a numeral. For
many proofs this matters little (we only care about propositional
equality), but for executable extraction it is fatal. Cubical type theory
solves this; some applications -- notably synthetic differential geometry
\cite{SchreiberShulman} -- additionally require modal axioms that
themselves are not yet known to be cubically computational. The trade-off
between axiomatic generality and definitional reduction is an active
research front.

\section{Companion Implementation}
\label{sec:haskell}

The companion Haskell library is available in the project's public
source repository (under \texttt{src/05-hott/}; see the
\emph{Univalent Correspondence} project). It provides:
\begin{itemize}
\item A GADT \texttt{Path a b} representing identity types with constructor
\texttt{Refl :: Path a a}.
\item Functions \texttt{sym}, \texttt{trans}, and \texttt{transport}
mirroring \Cref{lem:sym,lem:trans,def:transport}.
\item A type-level naturals module with peano-style \texttt{succ} and a
demonstration that 58 can be constructed and equal-typed.
\item A module sketching the circle as a HIT (necessarily approximated, as
Haskell lacks higher constructors) and the encode--decode method.
\item An Agda-style sketch in comments showing the analogous direct
type-theoretic syntax.
\end{itemize}
The \texttt{Main.hs} executable runs a small demonstration.

\section{Agda Sketch}
\label{sec:agda-sketch}

The most natural ambient language for HoTT is Agda or Cubical Agda. We
include here an Agda-flavoured sketch of the basic constructions; full
files for Cubical Agda live alongside the Haskell sources.

\begin{lstlisting}
-- Agda-style (comments only; the actual Agda lives elsewhere):
--
-- data Nat : Set where
--   zero : Nat
--   succ : Nat -> Nat
--
-- _==_ : {A : Set} -> A -> A -> Set
-- a == b = Path a b
--
-- refl : {A : Set} {a : A} -> a == a
-- refl = Refl
--
-- sym : {A : Set} {a b : A} -> a == b -> b == a
-- sym Refl = Refl
--
-- trans : {A : Set} {a b c : A}
--       -> a == b -> b == c -> a == c
-- trans Refl q = q
--
-- transport : {A : Set} (P : A -> Set)
--           -> {a b : A} -> a == b -> P a -> P b
-- transport P Refl x = x
--
-- ua : {A B : Set} -> A ~= B -> A == B
-- (postulated in Book HoTT, derived in Cubical)
\end{lstlisting}

\section{Conclusion and Future Work}
\label{sec:conclusion}

We have argued that Homotopy Type Theory provides the synthetic language
in which arithmetic, analysis, and algebraic topology share a single
foundational system. The natural number $58$ is the canonical
$\overline{58} : \N$; the universe-level identification of all NNOs is a
theorem; the transcendental constants $\pi$ and $e$ are unique-existence
centres in the higher inductive--inductive type $\R_c$. Univalence makes
the categorical and structural perspectives co-extensive with the
type-theoretic one.

The companion Haskell code provides a hands-on, executable mirror of the
key constructions. The next paper (Paper~06) develops the categorical /
structural perspective independently and shows that it merges with the
HoTT perspective via the internal language theorem; Paper~07 (the
synthesis) closes the loop.

\subsection{Open problems}

\begin{enumerate}
\item Formalize Lindemann--Weierstrass in Cubical Agda end-to-end.
\item Develop directed univalence sufficient to express the synthetic
$(\infty, 2)$-Yoneda lemma.
\item Formalize the Brown representability theorem inside synthetic
homotopy theory; this would clarify the relation to the spectra
constructions in classical algebraic topology.
\end{enumerate}

\begin{thebibliography}{99}

\bibitem{HoTTBook}
The Univalent Foundations Program.
\emph{Homotopy Type Theory: Univalent Foundations of Mathematics}.
Institute for Advanced Study, 2013. \texttt{arXiv:1308.0729}.

\bibitem{MartinLof84}
P.\ Martin-L\"of.
\emph{Intuitionistic Type Theory}. Bibliopolis, 1984.

\bibitem{LumsdaineGroupoid}
P.\ L.\ Lumsdaine. Weak omega-categories from intensional type theory.
\emph{Logical Methods in Computer Science}, 6(3:24):1--19, 2010.

\bibitem{vandenBergGarner}
B.\ van den Berg and R.\ Garner. Types are weak omega-groupoids.
\emph{Proc.\ London Math.\ Soc.}, 102(2):370--394, 2011.

\bibitem{LicataShulman}
D.\ R.\ Licata and M.\ Shulman. Calculating the fundamental group of the
circle in homotopy type theory.
In \emph{LICS 2013}, pp.\ 223--232.

\bibitem{LicataFinster}
D.\ R.\ Licata and E.\ Finster. Eilenberg--MacLane spaces in homotopy type
theory. In \emph{CSL-LICS 2014}, ACM, 2014.

\bibitem{Booij}
A.\ B.\ Booij. \emph{Analysis in univalent type theory}. PhD thesis,
University of Birmingham, 2020.

\bibitem{CCHM}
C.\ Cohen, T.\ Coquand, S.\ Huber, A.\ M\"ortberg.
Cubical type theory: a constructive interpretation of the univalence
axiom. In \emph{TYPES 2015}, LIPIcs 69, pp.\ 5:1--5:34.

\bibitem{CubicalAgda}
A.\ Vezzosi, A.\ M\"ortberg, A.\ Abel.
Cubical Agda: a dependently typed programming language with univalence and
higher inductive types.
\emph{Proc.\ ACM Program.\ Lang.}, 3(ICFP):87:1--87:29, 2019.

\bibitem{AwodeyWarren}
S.\ Awodey and M.\ A.\ Warren.
Homotopy theoretic models of identity types.
\emph{Math.\ Proc.\ Cambridge Philos.\ Soc.}, 146(1):45--55, 2009.

\bibitem{KapulkinLumsdaine}
C.\ Kapulkin and P.\ L.\ Lumsdaine.
The simplicial model of univalent foundations (after Voevodsky).
\emph{J.\ Eur.\ Math.\ Soc.}, 23(6):2071--2126, 2021.

\bibitem{ShulmanInftyTopos}
M.\ Shulman. All $(\infty,1)$-toposes have strict univalent universes.
\texttt{arXiv:1904.07004}, 2019.

\bibitem{ShulmanStrength}
M.\ Shulman. Higher inductive types and the formal Blakers--Massey theorem.
\emph{Proc.\ ACM Program.\ Lang.}, 4(POPL):41:1--41:30, 2020.

\bibitem{2LTT}
D.\ Annenkov, P.\ Capriotti, N.\ Kraus, C.\ Sattler.
Two-level type theory and applications.
\emph{Math.\ Structures Comput.\ Sci.}, 33(8):688--743, 2023.

\bibitem{SchreiberShulman}
U.\ Schreiber and M.\ Shulman. Quantum gauge field theory in cohesive
homotopy type theory. In \emph{QPL 2012}, EPTCS 158, pp.\ 109--126, 2014.

\bibitem{Riley}
M.\ Riley. \emph{A type theory with a tiny object}.
PhD thesis, Wesleyan University, 2022.

\bibitem{RiehlShulman}
E.\ Riehl and M.\ Shulman. A type theory for synthetic
$\infty$-categories. \emph{Higher Structures}, 1(1):147--224, 2017.

\bibitem{CisinskiNguyen}
D.-C.\ Cisinski, H.\ K.\ Nguyen.
The universal coCartesian fibration. \texttt{arXiv:2210.08945}, 2022.

\bibitem{AgdaUnimath}
The agda-unimath project. \emph{The agda-unimath library}.
\texttt{https://unimath.github.io/agda-unimath/}, accessed 2026.

\bibitem{SemisimplicialProblem}
N.\ Kraus. The general universal property of the propositional truncation.
\emph{TYPES 2014}, LIPIcs 39, pp.\ 111--145.

\bibitem{Voevodsky}
V.\ Voevodsky. The equivalence axiom and univalent models of type theory.
\texttt{arXiv:1402.5556}, 2014.

\bibitem{Wadler2015}
P.\ Wadler. Propositions as types.
\emph{Comm.\ ACM}, 58(12):75--84, 2015.

\bibitem{Lurie}
J.\ Lurie. \emph{Higher Topos Theory}.
Annals of Math.\ Studies 170, Princeton, 2009.

\end{thebibliography}

\end{document}
